%% 
%% Copyright 2019-2024 Elsevier Ltd
%% 
%% Version 2.4
%% 
%% This file is part of the 'CAS Bundle'.
%% --------------------------------------
%% 
%% It may be distributed under the conditions of the LaTeX Project Public
%% License, either version 1.2 of this license or (at your option) any
%% later version.  The latest version of this license is in
%%    http://www.latex-project.org/lppl.txt
%% and version 1.2 or later is part of all distributions of LaTeX
%% version 1999/12/01 or later.
%% 
%% The list of all files belonging to the 'CAS Bundle' is
%% given in the file `manifest.txt'.
%% 
%% Template article for cas-dc documentclass for 
%% double column output.

%\documentclass[a4paper,fleqn,longmktitle]{cas-dc}
\documentclass[a4paper,fleqn]{cas-dc}

%\usepackage[authoryear,longnamesfirst]{natbib}
%\usepackage[authoryear]{natbib}
\usepackage[numbers]{natbib}
\usepackage{array}

%%%Author definitions
\def\tsc#1{\csdef{#1}{\textsc{\lowercase{#1}}\xspace}}
\tsc{WGM}
\tsc{QE}
\tsc{EP}
\tsc{PMS}
\tsc{BEC}
\tsc{DE}
%%%Mine

\DeclareMathOperator*{\argmax}{arg\,max}
\DeclareMathOperator*{\argmin}{arg\,min}
\usepackage[ruled,vlined]{algorithm2e}
\hypersetup{
  colorlinks=False,
  citecolor=Violet,
  linkcolor=Red,
  urlcolor=Blue}

\usepackage{tabularx}
\usepackage{multirow}
\usepackage{makecell}
\usepackage{subcaption}
\usepackage{utfsym}
\usepackage{xcolor}
\usepackage{soul}
\usepackage{amsthm}
\usepackage{amsmath}
\newtheorem{assumption}{Assumption}[section]
\usepackage[switch]{lineno}

%%%%%%

\begin{document}
\emergencystretch 3em %%%exceed margin
\let\WriteBookmarks\relax
\def\floatpagepagefraction{1}
\def\textpagefraction{.001}
\shorttitle{MC production line simulation for human state estimation in task planning and allocation}
\shortauthors{Jintao Xue et~al.}

\title [mode = title]{Modular construction production line simulation for human state estimation in task planning and allocation}                 

\author[1]{Jintao Xue}[]
\author[1]{Xiao Li}[orcid=0000-0001-9702-4153]
\cormark[1]
\author[1]{Yuxiao Tian}[]
\author[1]{Hao Zheng}[]
\author[1]{Nianmin Zhang}[]
\affiliation[1]{organization={CI3 Lab, Department of Civil Engineering, The University of Hong Kong},
                city={Hong Kong SAR},
                }

\cortext[cor1]{Corresponding author: CI3 Lab, Department of Civil Engineering, Faculty of Engineering, HW6-07, Haking Wong Building, The University of Hong Kong, Pokfulam, Hong Kong, China.}


\begin{abstract}
In modular construction (MC) factories under Industry~5.0, human--robot collaboration (HRC) is increasingly adopted, with task planning and allocation (TPA) serving as a core decision layer. Yet most existing studies remain limited to small-scale settings or assume fully observable human states, while datasets that reflect real manufacturing process dynamics are scarce, hindering practical deployment. To address these gaps, we develop a factory-level MC production-line simulation platform in NVIDIA Isaac Sim, built from a replica of a real water-pipe production-part factory. The platform encodes the manufacturing workflow as a process-task graph that couples logistic and processing tasks, and decomposes each task into executable subtasks with explicit partner-agent dependencies among humans, gantries, machines, and robots. A rule-based multi-agent TPA controller then performs product sequencing, process-task planning, task allocation, and cross-agent progress updates, forming a closed-loop pipeline that synchronizes process execution with multi-camera observations and automatically generates process-consistent labels. On the resulting dataset, we study two human state estimation tasks: multi-view human identity recognition and working-human subtask recognition with completion status. Experiments show that the simulation stably provides reliable supervision over long-horizon episodes, and that models trained purely on synthetic data achieve high accuracy and strong generalization in simulation. These results indicate that MC production-line simulation is an effective data engine for perception-enabled, deployable TPA systems.
\end{abstract}

\begin{keywords}
Modular construction \sep Human-centric Manufacturing \sep Factory Simulation \sep Task planning and allocation \sep Vision-based state estimation
\end{keywords}

\begin{NoHyper}
\maketitle
\end{NoHyper}

%% ========================================================================
\section{Introduction} \label{sec:intro}
%% ========================================================================

Modular construction (MC) is a promising off-site construction paradigm. By shifting much on-site work to production-part factories, it improves productivity and quality while shortening project schedules~\cite{gibb1999off}.
Under Industry~5.0, MC factories increasingly adopt human--robot collaboration (HRC) following human-centric manufacturing~\cite{fu2024human}, where task planning and allocation (TPA) serves as a key decision layer~\cite{cheng2019task,banziger2020optimizing,xue2026safe}.
However, existing TPA research remains largely confined to small-scale settings, often a single robotic arm performing assembly, and offers limited support for line-scale MC production.
Moreover, the data required to train and evaluate perception-enabled TPA in MC factories remain scarce.
Despite substantial progress on TPA algorithms, integrating perception with TPA in large-scale settings with multiple humans, robots, and machines remains challenging and is still largely absent in the literature.

To address these gaps, we develop a factory-level simulation platform of an MC production line.
The platform is built in NVIDIA Isaac Sim replica of a real water-pipe production-part factory, including storage areas, gantry logistics, robotic welding stations, cutting and grooving machines, and workbenches. We encode the manufacturing workflow as a process-task graph with coupled logistic and processing tasks, and each task is further decomposed into executable subtasks with explicit partner-agent dependencies (human, gantry, machine, and robot). To keep the virtual factory running coherently, we implement a rule-based multi-agent TPA controller that performs product sequencing, process-task planning, human--robot task allocation, and cross-agent progress updates under coupling constraints. Under ideal perception, this controller drives a closed-loop data pipeline that synchronizes process execution with multi-camera observations and automatically generates process-consistent labels for training and evaluating human state estimation models, including human identity and working-human subtask states.
Experiments show that the rule-based controller can stably drive long-horizon production episodes and generate reliable process-consistent supervision, while models trained on the resulting synthetic dataset achieve high accuracy and good generalization in simulation.


In other words, the controller is assumed to know \emph{who} is present in each work zone and \emph{what} subtask each worker is currently performing.
This omniscient assumption simplifies algorithmic evaluation, yet it leaves a critical sensing gap between simulated decision policies and real factories, where such information must be recovered from cameras or other sensors~\cite{jahanmahin2022human,semeraro2023human}.
Without a perception module that can estimate human identity and subtask progress, TPA remains incomplete as a deployable closed-loop system.

A complementary obstacle is data scarcity.
Training industrial multi-view perception models requires large amounts of synchronized images whose labels are consistent with manufacturing process dynamics (task graphs, agent coupling, and temporal subtask transitions).
Collecting and annotating such data on a live MC line is costly, disruptive, and difficult to scale.
Photorealistic factory simulation offers a promising alternative~\cite{banziger2020optimizing}, but only if the simulated process itself is driven by the same TPA logic that would operate on the shop floor; otherwise, the generated visual streams may not reflect realistic agent interactions.

To address these gaps, this paper develops a photorealistic MC production-line simulation in which a rule-based TPA controller executes the end-to-end manufacturing workflow.
The closed-loop process continuously drives synchronized multi-camera observations and automatically emits ground-truth labels for human presence and subtask states, thereby serving as a scalable data engine for perception learning.
On the generated dataset, we study two complementary multi-view tasks:
Human identity recognition: For each calibrated camera, predict which humans appear in the view (multi-label classification over a fixed human vocabulary).
Human subtask recognition: For each working human, infer the current subtask class and a binary completion flag from multi-camera images conditioned on the allocated process task.

The main contributions are summarized as follows.
\begin{itemize}
    \item We develop a realistic MC production-line simulation platform addressing the current lack of data tailored to MC production TPA and perception learning.
    \item We couple rule-based TPA with a photorealistic MC line simulation so that process execution, multi-view imaging, and ground-truth labeling form a closed data-generation loop.
    \item We study two key perception tasks for line-level TPA: per-view human identity and subtask recognition, and show that models trained on simulator data achieve strong performance, enabling perception-driven HRC systems.
\end{itemize}

The remainder of this paper is organized as follows.
Section~\ref{sec:review} reviews related work and research gaps.
Section~\ref{sec:problem} formalizes the perception problems.
Section~\ref{sec:method} presents the simulation-driven data engine and learning pipelines.
Section~\ref{sec:exp} reports experiments, and Section~\ref{sec:conclusion} concludes.


%% ========================================================================
\section{Literature review} \label{sec:review}
%% ========================================================================

\subsection{Human-robot collaboration in MC production}
\label{sec:review_hrc}
Human--robot collaboration (HRC) is a key enabler for Industry~5.0, as humans and robots complement each other: humans contribute flexibility in complex and adaptive tasks, while robots provide efficiency in repetitive, hazardous, and physically demanding operations~\cite{xu2021industry, xueproduction}.
These systems also differ in interaction level, ranging from fully separated operation and human--robot coexistence with safety intervention, through assistive and cooperative modes with increasing workspace sharing, to simultaneous collaboration and largely autonomous execution~\cite{mukherjee2022survey,rodrigues2023multidimensional}.
In manufacturing, HRC research is commonly organized around four themes: (i)~collaborative workstation and cell design, including safety and ergonomics~\cite{tsarouchi2017human,wang2025design}; (ii)~task planning and allocation to support smooth collaboration and improve productivity, safety, and related performance metrics~\cite{zhao2025safety,zheng2023video,xue2026safe}; (iii)~perception and human behavior modeling, including human-aware state estimation~\cite{prunet2024optimization,jahanmahin2022human,semeraro2023human}; and (iv)~human--robot interaction and control~\cite{cherubini2016collaborative,li2023proactive}.
Within MC production specifically, the automation mix further expands to robotic arms, specialized robots, exoskeletons, and automated guided vehicles (AGVs), each supporting different logistics, processing, and assistance roles alongside human operators~\cite{fu2024human,malik2019generation,huysamen2018assessment,yang2021agv}.


\subsection{Task planning and allocation}
Task planning and allocation (TPA) is a core manufacturing decision layer that allocates operations to humans and automation assets under productivity, safety, and ergonomic constraints by determining what to execute, when to execute it, and which agent should perform each task~\cite{cheng2019task,banziger2020optimizing}.
Existing studies fall into two broad categories:
(i)~methods that focus on the TPA algorithm itself and treat perception as ideal or fully observable information, sometimes at relatively large problem scales~\cite{cai2023task,asadayoobi2023optimising,ye2024two,ren2023decision,yasari2025ergonomics,jia2025flexible,xue2026hierarchical}; and
(ii)~methods that incorporate perception, but only in small-scale settings, such as a single robot paired with a single machine or a narrow task type (mainly assembly)~\cite{lou2024human,wang2025transfer,merlo2023ergonomic,lee2022digital,zhao2025safety,zhang2022reinforcement}.
Many line-level formulations in the first category consume privileged state variables, such as agent locations, task progress, and completion flags, as if they were directly observed.
This practice is reasonable for algorithmic benchmarking, but it postpones the sensing problem that dominates real deployments.
Line-scale TPA research for MC production remains especially limited.
Although single-arm collaboration is relatively easy to prototype and validate, multi-agent, multi-machine production lines are costly and disruptive to test in live factories.
Consequently, work that jointly couples perception and TPA at line scale remains scarce, and scalable manufacturing simulation platforms are increasingly needed to develop and evaluate TPA algorithms under realistic process dynamics.

\subsection{Factory simulation and synthetic perception data}
\label{sec:review_sim}

Manufacturing simulation has long been used to evaluate schedules, layouts, and human--robot task allocation offline~\cite{banziger2020optimizing,ferjani2017simulation,pascual2024development}.
Workstation digital twins further combine physics-based rendering with ergonomics and safety analysis~\cite{wang2025design,lee2022digital}.
Recent reviews also emphasize synthetic data for manufacturing vision, yet most pipelines target defect detection or assembly perception at single-station scale~\cite{shafiee2026synthetic}.
Line-level simulators that couple multi-agent logistics with process-synchronized multi-camera labels remain scarce.

Recent Physical AI platforms, such as NVIDIA Isaac Sim, NuRec, SAGE, and Factory assembly benchmarks, enable photorealistic rendering and scalable GPU simulation for real-to-sim-to-train workflows~\cite{isaacsim2026,liang2018gpu,nurec2026,xia2026sage,narang2022factory}.
SAGE, for instance, generates simulator-ready 3D scenes from a user-specified embodied task through generator--critic refinement of semantic, visual, and physical validity~\cite{xia2026sage}.
However, these tools and public datasets still focus mainly on autonomous driving, warehousing, and manipulation.
NuRec and SAGE provide strong visual or layout control, but they do not offer process-controllable shop-floor equipment under a line-level TPA controller.
Process-consistent labeling for MC manufacturing HRC therefore remains an open need.

\subsection{Vision-based human state estimation in industrial settings}
\label{sec:review_vision}

Vision methods for human behavior modeling in manufacturing range from pose estimation and action recognition to intent prediction~\cite{jahanmahin2022human,semeraro2023human}.
Recent human-centric digital twin studies further recover operator status from RGB or RGB-D cameras and synchronize it into collaborative work-cell models.
Skeleton-based pipelines support safety distancing and motion or action recognition~\cite{li2022framework,choi2022integrated,wang2021digital}, while vision-based human digital twins estimate posture, intention, ergonomic risk, and fatigue-aware operations~\cite{fan2023vision,yao2024virtual,zhang2024enabling,chand2024vision}.
Wang et al.~\cite{wang2025design} construct a customized high-fidelity mesh human from a short RGB video and use transformers for real-time motion tracking and procedure recognition under occlusion.
These advances confirm that vision is a practical sensing channel for human-centric manufacturing.
However, most systems remain at single work-cell or station scale and mainly estimate local motion or procedure states, rather than the line-level observability required by TPA: multi-view human identity and process-coupled subtask completion.


\subsection{Research gaps}
\label{sec:gaps}

Drawing on the above review, three coupled gaps remain.
First, line-scale TPA for MC production with multiple humans, robots, and machines is rarely studied together with perception.
Second, many TPA methods assume omniscient human states (\emph{who} is present and \emph{what} subtask is ongoing), leaving a sensing gap for deployable closed-loop systems.
Third, process-consistent multi-view training data are scarce, and existing simulators seldom couple an executable TPA controller with automatic camera labeling.
This work addresses these gaps via a TPA-driven MC line simulation that serves as a data engine for multi-view human identity and subtask recognition.


%% ========================================================================
\section{Problem Formulation}\label{sec:problem}
%% ========================================================================

Consider an MC production line with a set of humans $\mathcal{H}=\{h_1,\ldots,h_{N_h}\}$, material-handling agents, and processing machines.
A rule-based TPA controller maintains ongoing task records and, for each working human $h\in\mathcal{H}$, exposes a process task $y^{\mathrm{task}}_h$ and a human-column subtask $y^{\mathrm{sub}}_h\in\mathcal{S}$, where $\mathcal{S}$ is a finite vocabulary of operational subtasks (e.g., travel to material, control gantry, wait).
A binary flag $y^{\mathrm{done}}_h\in\{0,1\}$ indicates whether the current human subtask has finished according to the process coupling rules.

A set of synchronized cameras $\mathcal{C}=\{c_1,\ldots,c_{N_c}\}$ observes the line.
At discrete time $t$, camera $c$ produces an RGB image $I_{c,t}$.
We study two complementary estimation problems that recover the human-centric observability required by TPA.

\paragraph{Task A: Human identity recognition.}
For each camera $c$ that participates in identity labeling, predict a multi-hot vector $z_{c,t}\in\{0,1\}^{N_h}$ whose $i$-th entry indicates whether human $h_i$ is visible in $I_{c,t}$.
Ground truth is obtained by testing each human's floor-plane position against a manually calibrated ground footprint polygon of camera $c$.

\paragraph{Task B: Human subtask recognition.}
For each working human $h$ at time $t$ (i.e., not free), estimate
\begin{equation}
\hat{y}^{\mathrm{sub}}_{h,t},\quad
\hat{y}^{\mathrm{done}}_{h,t}
=
f\!\left(
\{I_{c,t}\}_{c\in\mathcal{C}},\,
y^{\mathrm{task}}_{h,t}
\right),
\end{equation}
where $f$ is a multi-view network that fuses camera features with an embedding of the allocated process task.
The learning objective combines multi-class cross-entropy on $\mathcal{S}$ with binary cross-entropy on the done flag.

Together, Tasks~A and~B provide the missing perception interface between visual sensing and TPA: \emph{who} is in each view, and \emph{what} each working human is currently doing (and whether that subtask is complete).


%% ========================================================================
\section{Research method} \label{sec:method}
%% ========================================================================

\subsection{TPA-driven MC line simulation}
\label{sec:method_sim}

We reproduce a real MC production-part factory in a photorealistic simulator built on NVIDIA Isaac Sim / Isaac Lab.
The scene includes storage areas, gantries, robots, processing machines (e.g., welding, cutting, and grooving stations), workbenches, and multiple human workers.
Manufacturing products such as water-pipe assemblies progress through a process task gallery; each process task decomposes into coupled subtasks across human, gantry, machine, and robot columns.

A rule-based TPA controller executes the end-to-end workflow:
it allocates tasks to free humans, advances subtask rows according to coupling rules (e.g., human \texttt{control\_gantry} completes when the partner gantry reaches a carry/process state), and updates privileged progress records at every simulation step.
Because the same controller drives both logistics and machine operations, visual observations remain synchronized with process dynamics by construction.

\subsection{Multi-camera sensing and automatic labeling}
\label{sec:method_cam}

Fixed cameras are mounted at highrise overviews, the storage area, and key machine stations.
Each camera is configured with extrinsic pose $(\mathrm{eye},\mathrm{lookat})$ and RGB intrinsics.
For identity labeling, station-level cameras further define a manual ground footprint polygon $\mathrm{footprint}_c\subset\mathbb{R}^2$ on the factory floor.
At each logged step, human floor positions are tested by point-in-polygon queries; humans whose $(x,y)$ lie inside $\mathrm{footprint}_c$ form the multi-label ground truth for camera $c$.
Highrise overview cameras may capture RGB for context while being excluded from identity labeling when footprints are not calibrated.

Subtask ground truth is read directly from the TPA progress manager:
for each non-free human, the ongoing human-column subtask name and finished flag are serialized together with the allocated process task identifier.
Synchronized RGB frames from all cameras are saved alongside a compact environment record, yielding one JSONL sample per simulation step.

\subsection{Perception learning pipelines}
\label{sec:method_learn}

\paragraph{Data engine workflow.}
Collection runs with the rule-based controller under headless multi-camera rendering.
Each episode writes a directory of per-step camera images and a \texttt{meta.jsonl} file aligned with a unified sample template that jointly stores Task~A and Task~B fields.
Episodes are split into train/validation/test partitions at the episode level to avoid temporal leakage.

\paragraph{Model for Task~A.}
We train a single-view ResNet-18 encoder with a linear multi-label head over the human identity vocabulary.
The loss is binary cross-entropy with logits.
At evaluation, predictions are thresholded at $0.5$, and we report element-wise accuracy overall as well as conditioned on empty versus occupied frames.

\paragraph{Model for Task~B.}
A shared ResNet-18 backbone encodes every camera image; per-camera features are projected and concatenated with a learned embedding of the process task id.
Two heads predict the subtask class and the done probability.
The training loss is
\begin{equation}
\mathcal{L}
=
\mathrm{CE}\!\left(\hat{y}^{\mathrm{sub}}, y^{\mathrm{sub}}\right)
+
\mathrm{BCE}\!\left(\hat{y}^{\mathrm{done}}, y^{\mathrm{done}}\right).
\end{equation}


%% ========================================================================
\section{Experiments} \label{sec:exp}
%% ========================================================================

\subsection{Dataset and training setup}
\label{sec:exp_setup}

We collect six manufacturing episodes under rule-based TPA, totaling $26{,}360$ synchronized multi-view steps.
Each step stores RGB from ten cameras (two highrise overviews plus eight station/storage views) together with automatic labels.
Human identity uses a five-person vocabulary; station cameras with calibrated footprints participate in Task~A.
Human subtask recognition uses a nine-class vocabulary
\{\texttt{go\_to\_material},
\texttt{material\_on\_gantry},
\texttt{control\_gantry},
\texttt{material\_on\_robot},
\texttt{go\_to\_goal\_area},
\texttt{material\_on\_goal\_area},
\texttt{go\_to\_processing\_machine},
\texttt{control\_machine},
\texttt{wait}\},
plus a binary done target.
Episodes are split $4$/$1$/$1$ for train/validation/test.
Models are optimized with AdamW for $20$ epochs (learning rate $1\times10^{-4}$, weight decay $1\times10^{-5}$, image size $224$).

\subsection{Results}
\label{sec:exp_results}

Table~\ref{tab:results} summarizes validation performance of the best checkpoints.

\begin{table}[pos=t]
\caption{Perception performance on simulator-generated multi-view data.}
\label{tab:results}
\centering
\begin{tabular}{@{}llc@{}}
\toprule
Task & Metric & Value \\
\midrule
A: Identity & Element-wise acc. & $99.95\%$ \\
A: Identity & Occupied-frame elem. acc. & $99.89\%$ \\
A: Identity & Empty-frame elem. acc. & $100\%$ \\
B: Subtask & Subtask-class acc. (9-way) & $89.93\%$ \\
B: Subtask & Done-flag acc. & $87.05\%$ \\
\bottomrule
\end{tabular}
\end{table}

Task~A is nearly saturated: empty and occupied frames are both recognized with very high element-wise accuracy, indicating that footprint-consistent presence labeling is learnable from single-view RGB in this calibrated setting.
Task~B is substantially harder because subtasks are defined by process semantics and multi-agent coupling rather than appearance alone; nevertheless, the multi-view model recovers the correct subtask class for nearly $90\%$ of working-human samples and predicts completion with about $87\%$ accuracy.
These results suggest that TPA-driven simulation can supply process-consistent supervision sufficient to train the perception modules that line-level TPA presupposes.

\subsection{Discussion}
\label{sec:exp_disc}

Two observations follow.
First, identity recognition benefits from geometric footprint labels that are cheap to calibrate once camera poses are fixed; the visual classifier then amortizes that calibration across long episodes.
Second, subtask recognition still leaves headroom, especially for visually subtle transitions (e.g., wait versus control) and for done-flag prediction, which depends on partner-agent states that may be only partially visible.
Future work includes stronger temporal models, explicit cross-agent relational features, and sim-to-real transfer evaluation on the physical MC line.


%% ========================================================================
\section{Conclusion} \label{sec:conclusion}
%% ========================================================================

This paper addressed the perception gap in line-level TPA for MC production-part factories.
By coupling a rule-based TPA controller with a photorealistic multi-camera simulation of a real MC line, we constructed a scalable data engine that jointly generates process-consistent images and ground-truth labels.
On the resulting dataset, multi-view models achieved high accuracy for human identity recognition and promising accuracy for human subtask and completion estimation.
The findings indicate that TPA-driven MC line simulation is an effective pathway toward perception modules that make HRC decision stacks more complete and deployable.
Closing the loop by feeding estimated identity and subtask states back into online TPA, and validating sim-to-real transfer, remain important next steps.


\section*{CRediT authorship contribution statement}
\textbf{Jintao Xue}: Methodology, Software, Data curation, Writing – original draft. \textbf{Xiao Li}: Conceptualization, Supervision, Writing – review $\&$ editing. \textbf{Nianmin Zhang}: methodology analysis, experiment design.

\section*{Declaration of Competing Interest} The authors declare that they have no known competing financial interests or personal relationships that could have appeared to influence the work reported in this paper. 

\section*{Data availability} Data will be made available on request.

\section*{Acknowledgments} 
The work described in this paper is supported by grants from Technology Cooperation Funding Scheme (TCFS) (Ref No.GHP/321/22SZ), The University of Hong Kong (Ref No.109002002), and Innovation and Technology Fund (ITF) (Ref No. TP/041/24LP).


%% Loading bibliography style file
\bibliographystyle{cas-model2-names}

% Loading bibliography database
\bibliography{cas-refs}

\end{document}
